% APA 7 (§2.12): the reference list starts on a new page, heading centred and bold.
\FloatBarrier
\newpage
% The heading is a centred box, not a \section, so it never reached the contents
% page on its own. Register it by hand, after \newpage so the page number is right.
\phantomsection
\addcontentsline{toc}{section}{References}
{\centering\normalfont\bfseries\large References\par}
\nopagebreak

\begingroup
\sloppy
\setlength{\parindent}{0pt}
\everypar{\hangindent=0.5in\hangafter=1}

El Emam, K., Jonker, E., Arbuckle, L., \& Malin, B. (2011). A systematic review of re-identification
attacks on health data. \emph{PLoS ONE, 6}(12), e28071.
\url{https://doi.org/10.1371/journal.pone.0028071}\par

El Yagoubi, F., Badu-Marfo, G., \& Al Mallah, R. (2026). AgentLeak: A benchmark for internal-channel
privacy leakage in multi-agent LLM systems. \emph{IEEE Access, 14}, 94960--94978.
\url{https://doi.org/10.1109/ACCESS.2026.3704541}\par

HL7 International. (2024). \emph{SMART App Launch implementation guide} (Version 2.2.0).
\url{https://hl7.org/fhir/smart-app-launch/}\par

Jacob, D., Alghamdi, E., Hu, Z., Alomair, B., \& Wagner, D. (2025). \emph{Preventing prompt injection
with type-directed privilege separation} [Preprint]. arXiv.
\url{https://doi.org/10.48550/arXiv.2509.25926}\par

Johnson, A., Bulgarelli, L., Pollard, T., Horng, S., Celi, L. A., \& Mark, R. (2023). \emph{MIMIC-IV
Clinical Database Demo} (Version 2.2) [Data set]. PhysioNet.
\url{https://doi.org/10.13026/dp1f-ex47}\par

Johnson, A., Pollard, T., Badawi, O., \& Raffa, J. (2021). \emph{eICU Collaborative Research Database
Demo} (Version 2.0.1) [Data set]. PhysioNet. \url{https://doi.org/10.13026/4mxk-na84}\par

% APA 7 (§9.47): "Johnson, A." and "Johnson, A. E. W." are different authors and are
% ordered by their initials, so both "Johnson, A." entries precede this one. The
% initials are faithful to each source and are deliberately not harmonised.
Johnson, A. E. W., Bulgarelli, L., Shen, L., Gayles, A., Shammout, A., Horng, S., Pollard, T. J., Hao,
S., Moody, B., Gow, B., Lehman, L. H., Celi, L. A., \& Mark, R. G. (2023). MIMIC-IV, a freely
accessible electronic health record dataset. \emph{Scientific Data, 10}(1), Article 1.
\url{https://doi.org/10.1038/s41597-022-01899-x}\par

Juneja, G., Pasupulati, J. N. S., Albalak, A., Hua, W., \& Wang, W. Y. (2025). \emph{MAGPIE: A benchmark
for multi-agent contextual privacy evaluation} [Preprint]. arXiv.
\url{https://doi.org/10.48550/arXiv.2510.15186}\par

Kahn, M. G., Callahan, T. J., Barnard, J., Bauck, A. E., Brown, J., Davidson, B. N., Estiri, H., Goerg,
C., Holve, E., Johnson, S. G., Liaw, S.-T., Hamilton-Lopez, M., Meeker, D., Ong, T. C., Ryan, P.,
Shang, N., Weiskopf, N. G., Weng, C., Zozus, M. N., \& Schilling, L. (2016). A harmonized data quality
assessment terminology and framework for the secondary use of electronic health record data.
\emph{eGEMs, 4}(1), 18. \url{https://doi.org/10.13063/2327-9214.1244}\par

Mandel, J. C., Kreda, D. A., Mandl, K. D., Kohane, I. S., \& Ramoni, R. B. (2016). SMART on FHIR: A
standards-based, interoperable apps platform for electronic health records. \emph{Journal of the
American Medical Informatics Association, 23}(5), 899--908. \url{https://doi.org/10.1093/jamia/ocv189}\par

Neupane, S., Mittal, S., \& Rahimi, S. (2025). \emph{Towards a HIPAA compliant agentic AI system in
healthcare} [Preprint]. arXiv. \url{https://doi.org/10.48550/arXiv.2504.17669}\par

Pollard, T. J., Johnson, A. E. W., Raffa, J. D., Celi, L. A., Mark, R. G., \& Badawi, O. (2018). The
eICU Collaborative Research Database, a freely available multi-center database for critical care
research. \emph{Scientific Data, 5}(1), Article 180178.
\url{https://doi.org/10.1038/sdata.2018.178}\par

Privacy of Individually Identifiable Health Information, 45 C.F.R. \S\S 164.502(b), 164.514(b), 164.514(c),
164.514(d) (2013). \url{https://www.ecfr.gov/current/title-45/part-164/subpart-E}\par

Saltzer, J. H., \& Schroeder, M. D. (1975). The protection of information in computer systems.
\emph{Proceedings of the IEEE, 63}(9), 1278--1308. \url{https://doi.org/10.1109/PROC.1975.9939}\par

Shi, T., He, J., Wang, Z., Li, H., Wu, L., Guo, W., \& Song, D. (2025). \emph{Progent: Securing AI agents
with privilege control} [Preprint]. arXiv. \url{https://doi.org/10.48550/arXiv.2504.11703}\par

Sweeney, L. (2000). \emph{Simple demographics often identify people uniquely} (Data Privacy Working
Paper No. 3). Carnegie Mellon University.
\url{https://dataprivacylab.org/projects/identifiability/paper1.pdf}\par

U.S. Department of Health and Human Services, Office for Civil Rights. (2012). \emph{Guidance regarding
methods for de-identification of protected health information in accordance with the Health Insurance
Portability and Accountability Act (HIPAA) Privacy Rule}.
\url{https://www.hhs.gov/hipaa/for-professionals/special-topics/de-identification/index.html}\par

Walonoski, J., Kramer, M., Nichols, J., Quina, A., Moesel, C., Hall, D., Duffett, C., Dube, K.,
Gallagher, T., \& McLachlan, S. (2018). Synthea: An approach, method, and software mechanism for
generating synthetic patients and the synthetic electronic health care record. \emph{Journal of the
American Medical Informatics Association, 25}(3), 230--238.
\url{https://doi.org/10.1093/jamia/ocx079}\par

Wells, B. J., Chagin, K. M., Nowacki, A. S., \& Kattan, M. W. (2013). Strategies for handling missing
data in electronic health record derived data. \emph{eGEMs, 1}(3), 1035.
\url{https://doi.org/10.13063/2327-9214.1035}\par

Weiskopf, N. G., \& Weng, C. (2013). Methods and dimensions of electronic health record data quality
assessment: Enabling reuse for clinical research. \emph{Journal of the American Medical Informatics
Association, 20}(1), 144--151. \url{https://doi.org/10.1136/amiajnl-2011-000681}\par

\endgroup
